\documentclass[11pt,a4paper]{article}
\usepackage[utf8]{inputenc} % Not needed with XeLaTeX
\usepackage{fontspec}
\usepackage[margin=0.75in]{geometry}
\usepackage{xcolor}
\definecolor{darknavy}{RGB}{0,0,100}
\usepackage[colorlinks=true, urlcolor=darknavy, linkcolor=darknavy]{hyperref}
\usepackage{enumitem}
% \usepackage{fontawesome}
\usepackage{fontawesome5}
\usepackage{titlesec}
\usepackage{academicons}

% % --- Font: Source Sans Pro via Type1 (works with pdfLaTeX)
% \usepackage[default,scale=1.0]{sourcesanspro} % body + sans by default
% \renewcommand{\familydefault}{\sfdefault}      % make sans the default family
% \usepackage[scaled]{helvet}
\usepackage[default,scale=0.95]{opensans}

% (optional) math that pairs reasonably with sans body
\usepackage{newtxsf} % sans math; or use newtxtext/newtxmath for serif math

\titleformat{\section}{\Large\bfseries}{\thesection}{1em}{}[\titlerule]
\titlespacing{\section}{0pt}{12pt}{6pt}

\newcommand{\me}[1]{\textbf{#1}} % highlight own name

\usepackage{needspace}
\setlength{\emergencystretch}{2em}

\usepackage{comicneue}

\begin{document}

\begin{center}
  % {\Huge\textbf{Pavel Popov}}\\[0.5em]
  % {\comicneue\Huge\textbf{Pavel Popov}}\\[0.5em]
  {\fontspec{JetBrains Mono}\Huge Pavel Popov}\\[0.5em]
  {\fontspec{JetBrains Mono}\footnotesize +1 (470) 920-1876 \quad | \quad Atlanta, GA 30318} \quad | \quad {paavali.popov@gmail.com}\\
  {\small \href{https://github.com/paavalipopov}{Github: paavalipopov} \quad | \quad
  % \href{https://www.linkedin.com/in/pavel-popov-83b69a123/}{\faLinkedin\ Pavel Popov} \quad
  \href{https://www.linkedin.com/in/pavel-popov-83b69a123/}{LinkedIn: pavel-popov-83b69a123} \quad | \quad
  \href{https://scholar.google.com/citations?user=4alUD5UAAAAJ}{Google scholar: 4alUD5UAAAAJ}}
\end{center}

\section*{Summary}
{\small \noindent PhD candidate in Computer Science (expected December 2027) with a background in applied mathematics and physics, specializing in interpretable deep learning for large multimodal neuroimaging and time-series data.
Over 10 years of research experience, with 10 peer-reviewed publications (including \textit{IEEE Transactions on Medical Imaging}, \textit{NeuroImage}, and \textit{Physical Review Applied}), a book chapter, two open-source Python packages, and talks at US and international conferences. Seeking applied ML engineer/researcher roles.}

\section*{Skills}
\begin{itemize}
[leftmargin=*,noitemsep]
\item \textbf{Machine Learning:} Deep learning architecture design (incl. transformers), time-series analysis \& forecasting, causal inference, explainable AI, regression, decision trees \& ensembles, clustering.
\item \textbf{LLMs \& Generative AI:} Working with AI agents; familiar with prompt engineering and LLMOps.
\item \textbf{Programming \& Tools:} Python (PyTorch, Scikit-learn, Pandas, NumPy, SciPy, Matplotlib, Seaborn), MATLAB, SQL, mongoDB, Spark, Git, Docker, Linux/HPC (SLURM), Excel, PowerPoint.
\item \textbf{Data Analysis:} EDA, statistics, feature engineering, frequent pattern mining, data visualization.
\item \textbf{Domain Expertise:} Neuroimaging (fMRI), signal processing, solid state physics, spintronics.
\end{itemize}

\section*{Education}
\noindent\textbf{Georgia State University} \hfill Expected December 2027\\
PhD in Computer Science, Machine Learning \& Neuroimaging Analysis \hfill GPA 3.89\\
MS in Computer Science \hfill August 2026\\
{\small\textit{Coursework:} Advanced Machine Learning, Deep Learning, Advanced Data Mining, Big Data Programming}\\[0.3em]
\textbf{Moscow Institute of Physics and Technology} \hfill June 2020\\
MS in Applied Mathematics and Physics, Spintronics \hfill GPA 3.92\\[0.3em]
\textbf{Moscow Institute of Physics and Technology} \hfill June 2018\\
BS in Applied Mathematics and Physics, Solid State Physics \hfill GPA 3.81

\section*{Experience}
\noindent{\textbf{Center for Translational Research in Neuroimaging and Data Science}} \hfill January 2022 --- Now\\
Graduate Research Assistant \hfill \textit{Atlanta, USA}\\
[-2em]
\begin{itemize}[leftmargin=*,noitemsep]
    \item Developed a robust DL classifier for medical diagnostics of functional MRI scans with state-of-the-art performance and 10x speedup over existing benchmarks.
    \item Designed an interpretable deep learning model of dynamic effective brain connectivity that simulates multivariate time-series data, enabling the identification of novel biomarkers.
    \item Designed and trained deep models (incl. transformers) for large-scale imaging \& time-series data.
    \item Built frameworks for distributed DL training and deployment on multi-GPU HPC clusters (SLURM).
    \item Coordinated work of cross-functional teams to develop ML tools useful in clinical practice.
    \item Presented and translated technical findings for general audiences.
\end{itemize}

\needspace{5\baselineskip}
\noindent\textbf{Institute of Radioengineering and Electronics} \hfill May 2017 --- September 2020\\
Research Assistant \hfill \textit{Moscow, Russia}\\
[-2em]
\begin{itemize}[leftmargin=*,noitemsep]
    \item Modeled complex signal propagation in novel hardware architectures based on wave logic.
    \item Developed mathematical models for control of high-frequency signal dynamics, providing theoretical validation for prototype hardware controllers.
    \item Ran large-scale numerical simulations to validate theoretical models.
\end{itemize}

\section*{Open Source Projects}
\noindent\textbf{\href{https://pypi.org/project/ml4fmri/}{ml4fmri}: Python toolkit for benchmarking DL models for fMRI analysis}\\
[-2em]
\begin{itemize}[leftmargin=*,noitemsep]
    \item Unified 13 DL models for multivariate time series analysis under a standardized training API.
    \item Automated an all-model benchmarking pipeline for any input time series for reproducible research.
\end{itemize}

\noindent\textbf{\href{https://pypi.org/project/brainbow/}{brainbow} --- Neuroimaging visualization utility}\\
[-2em]
\begin{itemize}[leftmargin=*,noitemsep]
    \item Created a command-line tool for visualizing and exploring 3D/4D neuroimaging data formats.
\end{itemize}

\section*{Conferences \& Workshops}
\newcommand{\confattend}[3]{\noindent\hspace*{1.2em}\textbf{#1}, \emph{#2} \hfill #3\par\vspace{0.05em}}
{\small
\confattend{GSU Research Center Brain Health Synergy Symposium}{Atlanta, GA, USA}{Dec 2025}
\confattend{IEEE Workshop on Machine Learning for Signal Processing (MLSP 2025)}{Istanbul, Turkey}{Sep 2025}
\confattend{Organization for Human Brain Mapping (OHBM 2025)}{Brisbane, Australia}{Jun 2025}
\confattend{IEEE Brain Discovery and Neurotechnology Workshop}{Chicago, IL, USA}{Oct 2024}
\confattend{Organization for Human Brain Mapping (OHBM 2023)}{Montreal, Canada}{Jul 2023}
\confattend{Machine Learning for Science \& Engineering Workshop}{Seattle, WA, USA}{Jun 2023}
\confattend{Euro-Asian Symposium «Trends in Magnetism» (EASTMAG 2019)}{Ekaterinburg, Russia}{Sep 2019}
\confattend{International Congress on Ultrasonics (ICU 2019)}{Brugge, Belgium}{Sep 2019}
}

\section*{Publications}
{\small
\begin{enumerate}[leftmargin=*,itemsep=1pt,label=\arabic*.]

\item \me{P. Popov}, C. Morasso, G. Dolci, I. Boscolo Galazzo, G. Pearlson, D. Danks, V. Calhoun, G. Menegaz, and S. Plis, ``Message in a Bottleneck: Interpretable Deep Learning for Dynamic Effective Brain Connectivity,'' \textit{IEEE Transactions on Medical Imaging}, 2026.

\item Z. Iqbal, M.M. Rahman, Q. Zia, \me{P. Popov}, Z. Fu, V.D. Calhoun, and S. Plis, ``Interpretable Disorder Signatures: Probing Neural Latent Spaces for Schizophrenia, Alzheimer's, and Autism Stratification,'' \textit{Brain Sciences}, 2025.

\item \me{P. Popov}, C. Morasso, G. Dolci, I. Boscolo Galazzo, G. Pearlson, D. Danks, V. Calhoun, G. Menegaz, and S. Plis, ``DECIFRA: Deep Extraction of Causally Informed Features via Restricted Architecture,'' \textit{IEEE Int. Workshop on Machine Learning for Signal Processing (MLSP)}, 2025.

\item J. Van Duyn, \me{P. Popov}, A. Iraji, Z. Fu, V.D. Calhoun, S. Plis, and A. Fedorov, ``Boltzmann-Inspired Model for fMRI Time-Series Classification,'' \textit{Machine Learning for Health (ML4H)}, 2025.

\item \me{P. Popov}, U. Mahmood, Z. Fu, C. Yang, V. Calhoun, and S. Plis, ``A simple but tough-to-beat baseline for fMRI time-series classification,'' \textit{NeuroImage}, 2024.

\item \me{P.A. Popov}, A.R. Safin, A. Kirilyuk, S.A. Nikitov, I. Lisenkov, V. Tyberkevich, and A. Slavin, ``Voltage-Controlled Anisotropy and Current-Induced Magnetization Dynamics in Antiferromagnetic-Piezoelectric Layered Heterostructures,'' \textit{Physical Review Applied}, 2020.

\item A.R. Safin, S.A. Nikitov, A.I. Kirilyuk, D.V. Kalyabin, A.V. Sadovnikov, P.A. Stremoukhov, M.V. Logunov, and \me{P.A. Popov}, ``Excitation of Terahertz Magnons in Antiferromagnetic Nanostructures: Theory and Experiment,'' \textit{Journal of Experimental and Theoretical Physics}, 2020.

\item A.R. Safin, \me{P.A. Popov}, D.V. Kalyabin, and S.A. Nikitov, ``Synthesizer of Discrete Frequency Spectrum Based on an Antiferromagnetic Spintronic Oscillator,'' \textit{Technical Physics Letters}, 2020.

\item \me{P.A. Popov}, A.Yu. Sharaevskaya, E.N. Beginin, A.V. Sadovnikov, A.I. Stognij, D.V. Kalyabin, and S.A. Nikitov, ``Spin wave propagation in three-dimensional magnonic crystals and coupled structures,'' \textit{Journal of Magnetism and Magnetic Materials}, 2019.

\item S.A. Nikitov, A.I. Stognij, A.Yu. Sharaevskaya, A.V. Sadovnikov, \me{P.A. Popov}, D.V. Kalyabin, and E.N. Beginin, ``3D Magnonic Crystals,'' chapter in \textit{Three-Dimensional Magnonics: Layered, Micro- and Nanostructures}, CRC Press, 2019.

\item \me{P.A. Popov}, A.Yu. Sharaevskaya, D.V. Kalyabin, A.I. Stognii, E.N. Beginin, A.V. Sadovnikov, and S.A. Nikitov, ``Volume Magnetostatic Spin Waves in 3D Ferromagnetic Structures,'' \textit{Journal of Communications Technology and Electronics}, 2018.

\end{enumerate}
}

\end{document}
